\section{Overview}
\label{sec:overview}

\paragraph{Goal of the project.} Estimate the probability that a TPC track is matched to a BTOF hit
that it did not produce, by pairing tracks with BTOF hits from \emph{different} events of similar
topology (mixed events). This first step establishes the data access, software environment and
a QA baseline for the quantities the mixed-event study will use.

\paragraph{Main findings (test sample, \nEvtSel{} selected events).}
\begin{itemize}
\item Event selection keeps \nEvtSel{} of \nEvtAll{} events (\fracSel\%); the $|V_z|<70$\,cm cut
  dominates the loss (Table~\ref{tab:evtcut}). \nEvtTrig{} events carry the minimum-bias trigger 640051.
\item Mean refMult is \meanRefMult; mean number of BTOF hits stored per event is \meanNBTofHits,
  i.e.\ \hitsPerMatched{} BTOF hits per TOF-matched primary track (\perEvtTrkBeta{} per event).
  The number of matched tracks and of BTOF hits are strongly correlated (correlation factor \corrTofTrkHits).
\item Track selection (Table~\ref{tab:trkcut}): \nTrkEta{} primary tracks enter the QA sample
  (\perEvtTrkEta{} per event); \nTrkBeta{} of them have a BTOF match with valid $\beta$.
\item TOF matching efficiency (matched with $\beta>0$ / TPC tracks) is \effMid{} for $0.5<p_T<1.5$\,\GeVc,
  \effLow{} for $0.15<p_T<0.3$\,\GeVc, \effEtaCentral{} for $|\eta|<0.5$ and \effEtaOuter{} for $0.5<|\eta|<1$.
\item \mtwo{} peaks (for $0.4<p<1.0$\,\GeVc) are at \mTwoMeanPi, \mTwoMeanK{} and \mTwoMeanP{}\,GeV$^2/c^4$
  for $\pi$, K, p, consistent with PDG values within 0.005\,GeV$^2/c^4$.
\item \fracFlagTwo\% of matched tracks have \fcmd{btofMatchFlag}=2 (one hit, several track candidates).
\item Mixed events (Sec.~\ref{sec:me}): the module geometry fitted from production matches reproduces the
  production projection points to \geoFitRmsTest\,cm. With the picoDst helix, \meFracPMCellHiTest\% of production matches above
  1.5\,\GeVc{} are projected to the same cell. The random-match probability is \mePMENbpctTest\% (cell$\pm$1, production-like)
  and \mePMEpctTest\% (cell only), proportional to the number of BTOF hits in the event; it amounts to
  \meFakeMidNbTest\% of same-event matches for $0.5<p_T<2$\,\GeVc.
\item 1M-event sample (\meNRunsVthree{} runs, Sec.~\ref{sec:me1m}): $P_{\rm ME}$ = \mePMENbpctVthree\% (cell$\pm$1),
  $f_{\rm rand}$ = \meFakeMidNbVthree\% for $0.5<p_T<2$\,\GeVc; per-run $P_{\rm ME}$ divided by the occupancy stays within
  \mePMEoccMinVthree--\mePMEoccMaxVthree; the geometry refitted on 1M events agrees with the test-run geometry to \geoCmpMeanVthree\,cm.
\item $m^2$ (Sec.~\ref{sec:m2}): with the hit time taken from the picoDst matches and the production path-length formula,
  the same-event $m^2$ reproduces production (RMS difference \mtDmRmsVthree\,GeV$^2$/$c^4$). Mixed events overlaid on the
  track's own event, normalized per projected track and including hits without stored time, give random matches equal to
  \mtMENormFracCiVthree\% of the TOF matches in 0--5\% and \mtMENormFracCaVthree\% in 70--80\%. In 0--5\% at $0.8<p<1$\,\GeVc{}
  they are \mtFrPiMidCiVthree\% ($\pi$), \mtFrKMidCiVthree\% (K) and \mtFrPMidCiVthree\% (p) of the yield under the peaks.
  They explain only a small part of the SE yield between the proton and deuteron peaks (SE/ME = \mtSbRatioAllVthree).
\item Background origin (Sec.~\ref{sec:bkg}): the SE excess between the p and d peaks comes from neighbour-cell matches
  (\nbShareHI\% of it), not from dead cells, start time or pile-up. Production picks the largest-ToT hit among the projected
  and neighbour cells, so random hits can replace the track's own hit. Emulating this choice in picoDst (no ToT) brackets the
  in-time random matches at \mtMENormFracCiVeight--\mtMENormFracCiVseven\% of the TOF matches in 0--5\%. A late component flat in
  time, weakly dependent on centrality, has no time in the picoDst and cannot be mixed; this needs the MuDst.
\item MuDst (Sec.~\ref{sec:mudst}, \ovEvMtwo{} events of two runs): unmatched BTOF hits are \prUnmatchedFrac\% of all hits
  and much later than matched ones (\prUlate\% vs \prMlate\% above 26\,ns): slow, unreconstructed particles. Re-running
  \fcmd{StBTofMatchMaker}/\fcmd{StBTofCalibMaker} on MuDst reproduces all \clNSame{} production matches of the closure test
  (time of flight to \clDtofMed\,ps median). Main result, production makers re-run on the overlay of partner hits and
  partner tracks: random matches \ovtFrCi\% (0--5\%) and \ovtFrAll\% (0--80\%), SE/ME between p and d \ovtHiCi{}
  (\ovtHiMinCi--\ovtHiMaxCi{} per bin at 0.4--2\,GeV/$c$). Cross-checks: emulation \muFrProdCi\% and \muFrProdAll\%;
  without partner tracks \ovFrCi\%; the pool method (own hits removed) \poolFrCi\%, which is not a random-match fraction.
\item Large MuDst sample (Sec.~\ref{sec:big}, 102 files of 14 runs, \mubigEvMtwo{} events, method 1): random matches
  \mubigFrProdCi\% (0--5\%) and \mubigFrProdAll\% (0--80\%), SE/ME between p and d \mubigHiProdCi. In SE$-$ME the split Student-t
  line shape is best in \bigFitTtwoBestCi{} of \bigFitNbinsCi{} momentum bins (mean $\chi^2$/ndf \bigFitChiTtwoCi{} vs \bigFitChiGCi{}
  for Gaussians, 0--5\%); raw $\bar{\rm p}$/p = \bigRatPMin--\bigRatPMax{} at 0.4--2.5\,GeV/$c$. Method 2 (Sec.~\ref{sec:bigm2}) gives \ovbigFrCi\% and \ovbigFrAll\%, within \bigTwoOneFrDev\% of method 1, and
  yields within \bigTwoYieldAgree\% at 0.4--2\,GeV/$c$. The SE excess between p and d at low momentum is a fixed rate per track, independent of
  centrality and flat in time, so it is not a random match (Sec.~\ref{sec:lowp}).
\item Full simulation (Sec.~\ref{sec:sim}, HIJING+GEANT, \simEvGen{} events): \simFrDecA\% of the matches at 0.2--0.4\,GeV/$c$ are decay
  muons (they stay near the pion peak) and \simFrSecA\% secondaries, \simSecTray\% of which start in the BTOF trays. In the p--d window the
  correlated rate is \simCorrA{} per $10^4$ tracks, flat in centrality (data/simulation \simDataRatMinA--\simDataRatMaxA); at
  0.4--0.6\,GeV/$c$ the simulation is low (data/simulation \simDataRatMinB--\simDataRatMaxB; a simulated ToT closes part of it).
  Added to the data ME, the simulated component describes the data SE between p and d at 0.2--0.4\,GeV/$c$ (SE/(ME + sim.)
  \tmplFastRatSumCiA--\tmplFastRatSumPerA) but over-predicts the K--p valley and the region above the deuteron by 2--3$\times$, so it is
  not used as a correction.
\item Occupancy: trays \deadTrayList{} have no hits at all; \nDeadModule{} of 3840 modules
  (\fracDeadModule\%) are empty. No matched track points to an empty module.
\end{itemize}

\paragraph{Choices made (one line each; details where used).}
\begin{enumerate}
\item Production P24iy (latest of P21ic, P23id, P24iy in the file catalog) --- Sec.~\ref{sec:data}.
\item Input read through XRootD from the \texttt{local} storage replicas; the NFS catalog replicas
  checked were absent on disk --- Sec.~\ref{sec:data}.
\item STAR library SL24y run inside the SL7 container on the AlmaLinux~9 interactive node --- Sec.~\ref{sec:software}.
\item Minimum-bias trigger IDs 640001/11/21/31/41/51 --- Sec.~\ref{sec:evtcuts}.
\item Event cuts $|V_z|<70$\,cm, $V_r<2$\,cm; no $V_z^{\rm TPC}-V_z^{\rm VPD}$ cut --- Sec.~\ref{sec:evtcuts}.
\item Track cuts: primary, $p_T>0.15$\,\GeVc, nHitsFit$\ge15$, nHitsFit/nHitsMax$>0.52$, gDCA$<3$\,cm, $|\eta|<1$ --- Sec.~\ref{sec:trkcuts}.
\item TOF match: \fcmd{btofMatchFlag}$>0$ and $\beta>0$; primary momentum used for \mtwo --- Sec.~\ref{sec:tofdef}.
\item Approximate TOF $n\sigma$ with a fixed $1/\beta$ resolution of 0.013 (QA only) --- Sec.~\ref{sec:tofdef}.
\item BTOF module geometry fitted from production matches (no \fcmd{StBTofGeometry}) --- Sec.~\ref{sec:me}.
\item Projection with the picoDst global helix; matching at cell level, ``cell'' and ``cell$\pm$1'' --- Sec.~\ref{sec:me}.
\item Mixing within the same run, 14 $V_z$ $\times$ 9 refMult classes, 5 partner events --- Sec.~\ref{sec:me}.
\item Large sample: one job per run ($\ge$60k events), bad runs excluded, catalog replicas deduplicated --- Sec.~\ref{sec:me1m}.
\item Centrality from StRefMultCorr (2019 19.6\,GeV), pile-up events excluded from $m^2$ --- Sec.~\ref{sec:m2}.
\item Random $m^2$ from overlay mixed events (partner hits added to the track's own event), normalized per projected
  track, all hits included (untimed hits assumed to have the timed-hit time distribution) --- Sec.~\ref{sec:m2}.
\end{enumerate}
